\section{What the supplied refutation of ETH changes}
\label{sec:eth}

\subsection{The precise quantitative premise}

OpenAI's supplied Theorem~1.1, with source label
\texttt{thm:main}, is stronger than its abstract formulation
\cite[Theorem~1.1]{OpenAIETH}.
For each fixed integer \(k\geq2\), it counts the satisfying assignments of a
\(k\)-CNF with \(n\geq20\) declared variables and \(m\) input clauses in
\[
 (n+m)^{O_k(1)}
 2^{O_k(n\log\log n/\sqrt{\log n})}
\]
bit operations and \((n+m)^{O_k(1)}\) bits of space.
For each fixed \(D>0\), if \(m\leq Dn\), the running time improves to
\[
 (n+m)^{O_{k,D}(1)}2^{O_{k,D}(n/\sqrt{\log n})}.
\]
These are theorem premises under the convention requested in this task.
No Lean replay or independent verification of the complete counting
algorithm is claimed here.

The constants, polynomial exponents, and thresholds may depend on the
\emph{fixed} width \(k\). The statement gives no bound uniform in
\(k=k(n)\). For fixed \(k\), duplicate clauses can be removed and
tautologies discarded without changing the count; the number of distinct
non-tautological clauses is at most
\(\sum_{j=0}^{k}2^j\binom nj=n^{O_k(1)}\).
Reading a repetitively encoded input still incurs its polynomial
input-length cost. In this section, \emph{subexponential} means
\(2^{o(n)}\) in the parameter explicitly named; it does not mean the
sometimes-used class
\(\bigcap_{\epsilon>0}\operatorname{DTIME}(2^{n^\epsilon})\).

\subsection{CX-E01: which exponential-time hypotheses are refuted}
\label{finding:CX-E01}

\paragraph{Finding.}
The supplied theorem refutes ETH, \(\#\mathrm{ETH}\), fixed-width SETH,
the analogous counting SETH, their bounded-error randomized versions,
Gap-ETH, and NSETH. Existing conditional lower-bound theorems remain
valid implications, but their refuted premises no longer establish
hardness. Neither side of \(\mathrm P\) versus \(\mathrm{NP}\) is
established by these deductions.

\paragraph{Proof.}
For fixed \(k\), choose a constant \(C_k\) bounding the exponential
coefficient in the theorem and put
\[
 h_k(n)=C_k n\log\log n/\sqrt{\log n}.
\]
Since \(h_k(n)/n\longrightarrow0\), the exact counter runs in
\(2^{o(n)}(n+m)^{O_k(1)}\). Its answer is positive exactly when the
formula is satisfiable and zero exactly when it is unsatisfiable.
With \(k=3\), this contradicts the positive exponential rate required
by ETH and by \(\#\mathrm{ETH}\).

SETH asserts that for every \(\epsilon>0\) some fixed \(k\) has no
algorithm of time
\(2^{(1-\epsilon)n}(n+m)^{O(1)}\).
Fix \(\epsilon=1/2\). For \emph{each} fixed \(k\),
\(h_k(n)\leq n/2\) for all sufficiently large \(n\).
Finitely many smaller values are absorbed into a \(k\)-dependent
constant. Thus every fixed width admits the required algorithm with
the same exponential coefficient \(1/2\). The negation of SETH does
not require its polynomial factor or threshold to be independent of
\(k\). The same proof applies to its counting analogue.

A deterministic algorithm is a randomized algorithm with error zero,
so the corresponding bounded-error hypotheses are also contradicted.
On the Gap-ETH promise, return YES precisely when the exact count is
positive. This distinguishes a satisfiable 3-CNF from one in which every
assignment violates a specified positive fraction of clauses. With
the usual \(m=O(n)\) promise, the fixed-density bound is directly
\(2^{o(n)}\operatorname{poly}(n)\). No approximation algorithm is
needed. Finally the zero-count test is a deterministic, hence also
nondeterministic, subexponential algorithm for fixed-width CNF
unsatisfiability. It contradicts NSETH.
The corresponding hypothesis formulations and their direct use of the
counter also appear in the supplied preprint
\cite[\S1.1.1--\S1.1.3]{OpenAIETH}.

For the logical scope, the formulas \(H\Rightarrow B\) and \(\neg H\)
entail neither \(B\) nor \(\neg B\): assigning either truth value to
\(B\) satisfies both formulas. Thus deleting the premise of a
conditional lower bound is not a construction of a faster algorithm
for its target. This observation is not an independence theorem for
any particular complexity conjecture.
The displayed running-time benchmark is superpolynomial, since
\[
 \frac{n\log\log n/\sqrt{\log n}}{\log n}\longrightarrow\infty.
\]
It is therefore not a polynomial-time bound for SAT. An upper bound
also supplies no superpolynomial lower bound for SAT. These deductions
establish neither \(\mathrm P=\mathrm{NP}\) nor
\(\mathrm P\ne\mathrm{NP}\).
\hfill\(\square\)

\paragraph{Sparsification and its overhead.}
The full fixed-width statement already applies without a sparsification
step. If a separate procedure produces at most \(2^{\epsilon n}\)
fixed-width formulas on at most \(n\) variables, each with at most
\(D_\epsilon n\) clauses, in
\(2^{\epsilon n}\operatorname{poly}_{k,\epsilon}(n+m)\) time,
pad each formula to \(n\) declared variables by forcing each new
variable to zero. This preserves its count and gives density at most
\(D_\epsilon+1\).
The fixed-density counter then gives the combined bound
\[
 \operatorname{poly}_{k,\epsilon}(n+m)
 2^{\epsilon n+O_{k,D_\epsilon}(n/\sqrt{\log n})}.
\]
This follows by multiplying the per-instance cost by the number of
instances and charging their construction.
For decision, a disjunction-equivalent list suffices.
For counting by addition, the represented satisfying-assignment sets
must partition the original set, or a separately proved counting
reconstruction must be provided: overlapping lists cannot simply
have their counts summed.
For each fixed \(\epsilon\) the displayed exponent is
\(\epsilon n+o(n)\).
Substituting \(\epsilon=\epsilon(n)\) additionally requires control of
the constants and of the reduction's uniformity.
None of these extra premises is needed for the direct use of
Theorem~1.1 above.

\paragraph{Limits of the refutation.}
The theorem does not give a
\(2^{(1-\delta)n}\operatorname{poly}(n,m)\) algorithm for
\emph{unrestricted-width} CNF-SAT, because its width may grow with \(n\).
Nor does the refutation of SETH logically refute OVC, the 3SUM or APSP
conjectures, or \(\mathrm{FPT}\ne\mathrm{W[1]}\).
An implication from SETH to a weaker conjecture has the wrong direction
for that inference. Unconditional lower bounds, NP-hardness
reductions, and results based on independent assumptions retain their
previous status.

\subsection{CX-E02: a parameter-preserving transfer criterion}
\label{finding:CX-E02}

\paragraph{Finding.}
Suppose a problem on inputs of length \(L\), with parameter \(u\),
has a polynomial-time reduction to fixed-width \(k\)-CNF with
\(N\leq au+b\) variables and \(N+M\leq L^{O(1)}\) variables and clauses
in the standard explicit encoding, where
\(a,b,k\) are fixed constants. Then it is decidable in
\[
 L^{O(1)}
 2^{O(u\log\log u/\sqrt{\log u})}
\]
time and polynomial space.
A parsimonious reduction gives the same bound for its witness count.
If the formula additionally has \(M=O(N)\) clauses, the
\(\log\log u\) factor is absent. Polynomially many such calls and
polynomial-time answer processing preserve these bounds.

\paragraph{Proof.}
Construct the formula and run the supplied counter or its positivity
test. For \(N\leq au+b\), after changing constants to cover finitely
many small values,
\[
 N\log\log N/\sqrt{\log N}
 =O(u\log\log u/\sqrt{\log u}).
\]
A polynomial-time reduction produces at most polynomially many bits,
so the prefactor is \(L^{O(1)}\).
Parsimoniousness identifies the two counts exactly.
The fixed-density part of Theorem~1.1 gives the stated improvement
when \(M=O(N)\).
Polynomially many calls change only the polynomial factor; their
workspaces can be reused.
\hfill\(\square\)

A reduction with \(N=\Theta(u^c)\), \(c>1\), instead gives a benchmark
exponent of order \(u^c\log\log u/\sqrt{\log u}\), which is not \(o(u)\).
Ordinary polynomial-time completeness therefore does not suffice to
transfer subexponential time in the original parameter.
Similarly, a \(2^{o(v)}\) graph algorithm is not a bound
\(f(k)v^{O(1)}\) in solution size \(k\), and a formula with \(O(v)\)
variables does not provide a \(2^{o(t)}\operatorname{poly}(v)\)
algorithm in treewidth \(t\).

For a claimed polynomial-exponent improvement, such as a
subquadratic algorithm in \(L\), the actual reduction cost,
solver prefactor, number of calls, and answer-processing cost must
each be bounded within that target time.
The notation \(L^{O(1)}\) alone cannot certify this.
For fixed-parameter tractability in a parameter \(t\), a sufficient
separate reduction would have fixed clause width,
\(N\leq g(t)\), and \(M\leq h(t)L^C\), and run in
\(h(t)L^C\) time, for computable \(g,h\) and a constant \(C\).
Substitution into Theorem~1.1 then gives
\(F(t)L^{O(1)}\) for a computable \(F\): every factor involving
the bounded variable count is absorbed into \(F(t)\).
The theorem does not itself construct such a reduction.

The supplied Corollary~1.3 already gives exact subexponential counting
of independent sets and fixed-\(q\) graph colorings in the vertex
parameter, and of perfect matchings in \(v+e\)
\cite[Corollary~1.3]{OpenAIETH}.
Bounded-set-size Set Splitting has a direct encoding with one bit per
universe element and, for each set, an all-positive and an all-negative
clause. These are actual algorithmic transfers.
Deleting a SETH premise alone supplies no subquadratic algorithm for
Edit Distance, LCS, or arbitrary Orthogonal Vectors.

\paragraph{A further explicit parameter check: counting cliques.}
For a graph with \(v\) vertices, the number of \(k\)-cliques can be
computed in
\(\operatorname{poly}(v+e)
 2^{O(v\log\log v/\sqrt{\log v})}\),
uniformly for \(0\leq k\leq v\).
Use a selection bit for every vertex, and for each nonedge
\(\{a,b\}\) impose \(\neg x_a\vee\neg x_b\).
Add a Boolean circuit requiring exactly \(k\) selected vertices.
A balanced tree of binary additions computes the population count
with \(O(v)\) gates: level \(j\) has \(O(v/2^j)\) additions on
\(O(j)\) bits, and
\(\sum_{j\geq1}vj/2^j=O(v)\).
Comparison with \(k\) adds \(O(\log v)\) gates.
The gate encoding in CX-E03 gives \(O(v)\) variables and \(O(v^2)\)
clauses, with a unique extension for each clique.
Theorem~1.1 proves the bound. This does not assert fixed-parameter
tractability in \(k\).
\hfill\(\square\)

\subsection{CX-E03: the Circuit-SAT saving and its exact parameter}
\label{finding:CX-E03}

\paragraph{Finding.}
An arbitrary Boolean circuit with \(n\) inputs, \(s\) noninput gates
of fan-in at most two, and unrestricted fan-out has a deterministic
satisfiability and exact accepting-input-count algorithm of time
\[
 (n+s)^{O(1)}
 2^{O((n+s)/\sqrt{\log(n+s)})}
\]
and polynomial space.
For \(s=O(n)\) this is \(2^{o(n)}\).
This bound does not meet the all-polynomial-size Circuit-SAT premise
of Williams's Theorem~1.1 \cite[Theorem~1.1]{Williams2013}.

\paragraph{Proof.}
Introduce one variable for each input and each gate.
A two-input gate has eight assignments to its inputs and output,
exactly four inconsistent with its truth function.
Forbid each inconsistent assignment by a clause of width at most
three; constants and unary gates need no more clauses.
Add a unit clause setting the final output to one.
The resulting formula has \(N=n+s\) variables and
\(M\leq4s+1\leq5N\) clauses.
Every input assignment extends to exactly one assignment to the gate
variables, by evaluation in topological order.
Thus this reduction preserves the number of accepting inputs.
Apply the fixed-density part of Theorem~1.1 with \(k=3,D=5\).
\hfill\(\square\)

Precisely, Williams's Theorem~1.1 assumes a superpolynomial function
\(q(n)\) such that Circuit-SAT on \(n\) inputs and \(n^a\) gates
has a deterministic or co-nondeterministic algorithm of time
\[
 \frac{2^n\operatorname{poly}(n^a)}{q(n)}
 \qquad\text{for every fixed }a.
\]
Its conclusion is \(\mathrm{NEXP}\not\subseteq\mathrm{P/poly}\).
For \(a>1\), substituting \(s=n^a\) into the proved benchmark gives an
exponent of order \(n^a/\sqrt{\log n}>n\).
Consequently this reduction does not supply that premise.
This is a diagnosis of the particular attempted inference, not an
impossibility theorem for other approaches to stronger lower bounds.

\subsection{CX-E04: the quantitative Boolean circuit lower bound}
\label{finding:CX-E04}

\paragraph{Finding.}
Under the task's theorem convention there exist a language
\(L\in\mathrm E^{\mathrm{NP}}\) and \(c>0\) such that
\[
 C_L(n)>c\,n\sqrt{\log n}
 \quad\text{for infinitely many }n,
\]
where \(C_L\) counts arbitrary fan-in-two Boolean gates.
In particular,
\(\mathrm E^{\mathrm{NP}}\not\subseteq\mathrm{SIZE}(O(n))\).
The quantitative statement is already
\cite[Corollary~1.4]{OpenAIETH}; it is not attributed as a new
circuit lower-bound discovery to this run.
The qualitative implication is
\cite[Theorem~6.1]{Williams2013}.
The finding here is the exact transfer and its scope.

\paragraph{Proof of the transfer.}
The decision algorithm required by the supplied Corollary~1.4 is
obtained from the fixed-density part of Theorem~1.1 by testing the
count for positivity.
Its premise is therefore satisfied, independently of both GRH
sources. Apply that corollary.
Since \(\sqrt{\log n}\longrightarrow\infty\), its language cannot
have a circuit family of eventual size \(O(n)\).
\hfill\(\square\)

\paragraph{Check of the cited implication proof.}
The complete proof in the supplied
\nolinkurl{quantitative-circuits.tex}, inserted in Appendix~C.1, was
read and its quantitative parameters checked.
It uses the unary nondeterministic time hierarchy in
\cite[Theorem~2.1]{Williams2013} and the local reduction
\cite[Theorem~1]{JMV2013}.
The latter gives a clause decoder with
\(q=n+O(\log n)\) index bits and \(O(n)\) output bits, each depending
on only constantly many index bits.
Enumerating inputs of bounded Hamming weight recovers supports and
truth tables, hence an \(O(n)\)-gate decoder in polynomial time.
The bits of the first accepting assignment, indexed by padded
\(dn\)-bit variable names, form a language in
\(\mathrm E^{\mathrm{NP}}\).
A hypothetical witness circuit of size
\[
 s_n=\left\lfloor\eta n
          \left\lfloor\sqrt{\log_2 n}\right\rfloor\right\rfloor
\]
is guessed. Three copies together with the decoder test whether a
clause is violated.
The gate encoding produces
\[
 U_n=(3\eta+o(1))n\sqrt{\log_2 n}
\]
variables and at most \(5U_n\) clauses.
For a solver bound \(V^b2^{AV/\sqrt{\log_2 V}}\), its cost is
\(2^{(3A\eta+o(1))n}\).
Choosing \(\eta<1/(12A)\) leaves enough margin, including polynomial
overhead, to contradict the chosen
\(\mathrm{NTIME}(2^n)\setminus\mathrm{NTIME}(2^{n/2})\)
language.
No gap was found in this implication proof.
This is not a claimed reproof of the counter, the hierarchy theorem,
or the local-reduction theorem.

The result is a bound in \(\mathrm E^{\mathrm{NP}}\), not in NP;
it concerns nearly linear Boolean circuits, not all polynomial-size
circuits or arithmetic circuits for the permanent.
The displayed existential constant \(c\) and infinitely-often
quantifier should not be strengthened to
\(\mathrm E^{\mathrm{NP}}\not\subseteq
  \mathrm{SIZE}(O(n\sqrt{\log n}))\).

\subsection{CX-E05: a quantitative sparse permanent algorithm}
\label{finding:CX-E05}

\paragraph{Finding.}
For \(A\in\{0,1\}^{r\times r}\) with \(e\) nonzero entries, its
integer permanent can be computed in
\[
 \operatorname{poly}(r+e)
 2^{O((r+e)/\sqrt{\log(r+e)})}
\]
time and polynomial space.
When \(e=O(r)\), the bound is
\(\operatorname{poly}(r)2^{O(r/\sqrt{\log r})}\).
This makes quantitative the matching consequence of the supplied
Corollary~1.3 by using its sharper fixed-density Theorem~1.1.

\paragraph{Proof.}
Interpret \(A\) as a bipartite graph with \(2r\) vertices and \(e\)
edges, with one selection bit for each edge.
At a vertex of degree \(d\), maintain bits \(a_i,b_i\) saying that
among the first \(i\) incident edges, respectively, at least one and
at least two have been selected. Initialize \(a_0=b_0=0\), and set
\[
 a_i=a_{i-1}\vee x_i,\qquad
 b_i=b_{i-1}\vee(a_{i-1}\wedge x_i).
\]
Induction on \(i\) proves these meanings.
Exactly one incident edge is selected if and only if
\(a_d\wedge\neg b_d=1\).
This requires \(O(d+1)\) gates.
Summing degrees, and conjoining the tests, gives \(O(r+e)\) gates;
an isolated vertex may instead be rejected immediately.

Apply the parsimonious encoding of CX-E03.
It has \(N=O(r+e)\) variables and \(M=O(N)\) clauses.
Every perfect matching gives one satisfying assignment with uniquely
determined gate bits, and every satisfying assignment selects a
perfect matching.
The number of these matchings is
\[
 \operatorname{per}(A)=
 \sum_{\sigma\in S_r}\prod_{i=1}^r A_{i,\sigma(i)}.
\]
The fixed-density bound proves the assertion.
\hfill\(\square\)

For dense matrices this encoding has \(\Theta(r^2)\) variables; its
bound does not give \(2^{o(r)}\) time in matrix dimension.
Nor does it construct an arithmetic circuit for the formal
polynomial \(\operatorname{per}_r(X)\).
The algorithm evaluates an integer count at a supplied 0--1 matrix.

\subsection{CX-E06: VP, VNP, and the surviving GCT gates}
\label{finding:CX-E06}

\paragraph{Finding and exact logical scope.}
Neither \(\mathrm{VP}=\mathrm{VNP}\) nor
\(\mathrm{VP}\ne\mathrm{VNP}\), and neither a polynomial arithmetic
circuit nor a superpolynomial arithmetic lower bound for the
permanent, is established by CX-E01--CX-E05.
The deductions do not discharge the border-hardness or black-box
identity-testing assumptions in GCT V Theorems~1.7, 2.4, 5.11,
and~5.14. This is a statement about the premises of these particular
implications, not a formal independence theorem.

\paragraph{Dictionary and proof.}
Over a fixed field, VP consists of polynomial-degree families with
polynomial-size arithmetic circuits. VNP permits Boolean-cube sums
of such polynomial families.
For characteristic other than two, the permanent is VNP-complete
under polynomial projections
\cite[Theorem~2, p.254]{Valiant1979}.
The determinant is complete for VBP, the weakly skew or branching
program subclass; it is not thereby complete for all of VP.
The corresponding orbit-closure formulation uses the border of
that model
\cite[Definitions~2.9, 2.25; Theorems~2.20, 2.29, 7.3]{Burgisser2024}.

The supplied counter accepts a finite bit-string description of a
fixed-width CNF and returns an integer.
Its conclusion does not construct, from an arithmetic verifier
\(g_n(X,Y)\), an arithmetic circuit for the symbolic polynomial
\[
 \sum_{Y\in\{0,1\}^{p(n)}}g_n(X,Y).
\]
CX-E05 also evaluates only a numerical specialization.
The lower bound in CX-E04 has a language in
\(\mathrm E^{\mathrm{NP}}\) as its target and Boolean gates as its
model. It has neither the family nor the model required for an
arithmetic permanent lower bound.
Thus these deductions do not furnish the premises of the cited
algebraic implications.

There is a separate quantitative obstruction to a tempting inference.
GCT V Theorem~1.7, source label \texttt{tintronnldetnew}, assumes
that for some fixed \(\epsilon>0\) the permanent cannot be approximated
infinitesimally by affine-linear symbolic determinants of size
at most \(2^{r^\epsilon}\).
Even if a new argument provided an upper bound
\(2^{O(r/\sqrt{\log r})}\) in that same determinant or border model,
it would not refute
\(2^{r^\epsilon}\) hardness for \(0<\epsilon<1\), because
\[
 \frac{r/\sqrt{\log r}}{r^\epsilon}
 =\frac{r^{1-\epsilon}}{\sqrt{\log r}}\longrightarrow\infty.
\]
No such symbolic conversion has been established here in the first
place.
The remaining gate is a substantive arithmetic conversion or
obstruction construction, not an erased GRH assumption.
\hfill\(\square\)

The GCT assumptions just mentioned are exactly located in the supplied
source as \texttt{tintronnldetnew}, \texttt{trussell2},
\texttt{tnnlexpfull}, and \texttt{tequisdit}, respectively.
Improving the complexity of finding a normalization system is also
distinct from proving the orbit-closure noncontainment required by a
permanent lower bound.

